\section{降维与升维的对偶：在合适的维度做合适的事}

前几篇把降维写成``压缩信息、恢复本质''。但机器学习还有一条相反的路线：核方法把数据映射到更高维（甚至无穷维），深度学习逐层升维再降维。这两条路线看似矛盾，实则互补：低维适合看清形态，高维适合精细交互。理解这个对偶，是理解什么时候该降维、什么时候该升维的关键。

\subsection{先感受困难：低维的清晰与高维的精细不可兼得}

考虑一个二分类问题：两个同心圆，内圈一类、外圈一类。在原始二维空间中，任何直线都无法分开它们；但映射到三维 $(x_1,x_2)\mapsto(x_1,x_2,x_1^2+x_2^2)$，数据变成可分的平面——高维让结构散开。

反过来，一个高维数据集（如 $1000$ 维基因表达）直接观察无法看出任何模式；PCA 降到二维，聚类结构立刻显现——低维让形态可见。

两个方向各有代价：升维可能引入维度诅咒（距离失真、样本稀疏），降维必然损失信息（不可逆、可能丢失任务相关结构）。选择哪个方向，取决于任务：理解数据形态，降维；精细计算结构，升维。

\subsection{降维：低维看清形态}

降维的目标是``恢复世界的真实维度''。流形假设（本单元第四篇）说：高维观测躺在低维流形上，降维是逼近这个流形。低维表示的价值在于：

\begin{itemize}
\tightlist
\item
  \textbf{可视化}：二维/三维投影让人直接看见聚类、异常、决策边界；
\item
  \textbf{结构发现}：PCA 的主方向、流形学习的测地距离，揭示数据生成的低维机制；
\item
  \textbf{去噪}：低秩近似滤除高维噪声，保留本质结构。
\end{itemize}

Johnson--Lindenstrauss 引理给降维一个理论保证：$n$ 个点在 $\R^d$ 中，可以随机投影到 $\R^k$，$k=O(\varepsilon^{-2}\log n)$，保持距离结构。投影维度只依赖点数，不依赖原始维度——高维数据的形态可以在低维忠实呈现，只要关心的是距离关系。

但降维是信息有损的：不可逆、可能丢失任务相关结构。t-SNE 的二维图可以看清聚类，但不能用来做分类；PCA 的前两个主成分可能丢弃判别方向。

\subsection{升维：高维精细交互}

升维的目标不是``看清''，而是``计算''。Cover 定理（1965）给出升维的数学依据：非线性可分的数据，在高维空间更可能线性可分。模式识别中``把数据映射到高维''的有效性，根源在这里。

核方法是升维的极致：不显式构造高维映射 $\phi(x)$，只用核函数 $K(x,x')=\langle\phi(x),\phi(x')\rangle$ 在高维空间做内积。SVM、核回归、Gaussian process 都在这个框架下工作——计算在高维（甚至无穷维），但不用显式存储高维坐标。

深度学习的架构也遵循这个对偶：卷积层、注意力层逐层升维（通道数增加、序列长度不变），最后全局池化或分类头降维（到类别数）。中间层在高维做``精细交互''（特征提取、信息混合），输出层在低维做``决策''（分类、回归）。

但升维的``精细''有条件：只在有结构的高维成立。流形结构（数据躺在低维流形上）、稀疏结构（大部分坐标为零）、低秩结构（矩阵可分解为低秩因子）让高维计算有效；无结构的高维（白噪声、加密数据）没有优势——距离失去意义、样本稀疏、计算昂贵。

\subsection{对偶的统一：计算在高维，理解在低维}

降维与升维不是二选一，而是同一过程的两个阶段：

\begin{enumerate}
\def\labelenumi{\arabic{enumi}.}
\tightlist
\item
  \textbf{计算在高维}：利用 Cover 定理（线性可分性）、表示定理（精确优化）、核技巧（隐式升维），在原始或映射后的高维空间做梯度下降、内积计算、特征提取；
\item
  \textbf{理解在低维}：利用 JL 引理（距离保持）、流形假设（结构恢复）、可视化（人类认知），把高维结果投影到低维，看清聚类、异常、决策边界。
\end{enumerate}

这个对偶在书里反复出现：

\begin{itemize}
\tightlist
\item
  \textbf{核方法}：高维计算（核函数内积），低维理解（支持向量、决策边界可视化）；
\item
  \textbf{深度学习}：中间层高维特征（精细交互），输出层低维决策（分类、回归）；
\item
  \textbf{表示定理}：无穷维优化问题的解落在有限维样本张成的子空间——计算在无穷维，解在低维；
\item
  \textbf{流形学习}：高维观测（精细结构），低维嵌入（形态可见）。
\end{itemize}

选择维度的标准不是``越高越好''或``越低越好''，而是``在合适的维度做合适的事''：需要精细计算时，升维；需要理解形态时，降维。

\subsection{边界清单}

\begin{itemize}
\tightlist
\item
  降维不可逆：低维表示不能完美重建高维数据，任务相关结构可能丢失；
\item
  升维有条件：只在有结构的高维（流形、稀疏、低秩）有效，无结构高维没有优势；
\item
  JL 引理保距离不保角度：随机投影保持距离结构，但不保证角度、面积等几何量；
\item
  Cover 定理是存在性结果：高维空间存在线性可分超平面，但不保证找到它（优化可能陷入局部极小）。
\end{itemize}

\subsection{落点}

降维与升维的对偶是机器学习的一个深层结构：低维适合看清形态（可视化、结构发现、去噪），高维适合精细交互（计算、优化、线性化）。理解这个对偶，就是理解什么时候该降维（理解数据）、什么时候该升维（精细计算）、以及什么时候两者结合（核方法、深度学习）。这是从``会用降维/升维工具''到``理解维度选择''的关键一步。
